% (User input window)
% 2026-02-15 14:39:13 (Asia/Singapore)

\begin{remark}[Moment Kernel Family for Endpoint Second-Order Audit]\label{rem:pom-diagonal-rate-endpoint-three-kernel-principle}
Theorem \ref{thm:pom-diagonal-rate-small-distortion-hellinger-endpoint} shows that the leading-order non-trivial constant at the small distortion endpoint \(\delta\downarrow 0\) is uniquely determined by \(A_{1/2}(w)^2-1\) (equivalently the half-order partition \(S_{1/2}\));
while Theorem \ref{thm:pom-diagonal-rate-zero-rate-endpoint-quadratic-curvature} and Corollary \ref{cor:pom-diagonal-rate-folded-curvature-auditable} show that the quadratic lift-off curvature constant at the zero rate endpoint \(\delta\uparrow\delta_0\) is determined only by \((p_2,p_3)\) (equivalently \((S_2,S_3)\)).
Therefore, if we seek only second-order precision characterization at both endpoints, we need to simultaneously introduce three moment kernel calibers: \(q=\tfrac12\) and \(q=2,3\).
\end{remark}

\begin{theorem}[Finite Integer Collision Moments Cannot Replace \texorpdfstring{$q=\tfrac12$}{q=1/2} Endpoint Constant]\label{thm:pom-finite-integer-moments-cannot-replace-half}
Given any fixed integer \(Q\ge 2\), there exist two families of fiber spectra \(\{d_m\}\), \(\{\widetilde d_m\}\) (which can be taken on the same \(X_m\), and satisfy \(\sum_x d_m(x)=\sum_x \widetilde d_m(x)=2^m\)), such that for all sufficiently large \(m\) we have
\[
\sum_{x\in X_m} d_m(x)^q=\sum_{x\in X_m}\widetilde d_m(x)^q\qquad(\forall q\in\{1,2,\dots,Q\}),
\]
but
\[
\sum_{x\in X_m} d_m(x)^{1/2}\ \ne\ \sum_{x\in X_m}\widetilde d_m(x)^{1/2}.
\]
Let the corresponding pushforward distributions \(w_m=d_m/2^m\), \(\widetilde w_m=\widetilde d_m/2^m\).
Then the finite integer-order collision moments \(p_q(\cdot)\) (\(q\le Q\)) of \(w_m,\widetilde w_m\) are completely identical, while the small distortion endpoint constants \(A_{1/2}(w_m)^2-1\) and \(A_{1/2}(\widetilde w_m)^2-1\) are strictly different;
thus the leading-order expansion of \(R_{w_m}(\delta)\) as \(\delta\downarrow 0\) (Theorem \ref{thm:pom-diagonal-rate-small-distortion-hellinger-endpoint}) cannot be determined by any integer moment audit mechanism of fixed order in the uniform sense as \(m\to\infty\).
\end{theorem}
\begin{proof}
First take a pair of distinct non-negative integer multisets \(A=\{a_i\}_{i=1}^n\), \(B=\{b_i\}_{i=1}^n\) such that the power sum matching
\[
\sum_{i=1}^n a_i^q=\sum_{i=1}^n b_i^q\qquad(q=0,1,\dots,Q)
\]
holds (for example, by Prouhet-Tarry-Escott construction we can give such equal-power-sum partitions for any \(Q\)).
Since \(A\neq B\), there exists a minimal integer \(q_0\ge Q+1\) such that \(\sum_i a_i^{q_0}\ne \sum_i b_i^{q_0}\).
\par
For any integer shift \(c\ge 1\), let \(A^{(c)}:=\{a_i+c\}\), \(B^{(c)}:=\{b_i+c\}\).
By binomial expansion,
\(
\sum_i (a_i+c)^q-\sum_i (b_i+c)^q
\)
can be expressed as \(\sum_{j=0}^q \binom{q}{j}c^{q-j}\bigl(\sum_i a_i^j-\sum_i b_i^j\bigr)\),
so for all \(q\le Q\) we still have \(\sum_i (a_i+c)^q=\sum_i (b_i+c)^q\).
On the other hand, let
\(
G(c):=\sum_i \sqrt{a_i+c}-\sum_i \sqrt{b_i+c}
\).
Using the binomial series
\(
\sqrt{c+t}=c^{1/2}\sum_{k\ge 0}\binom{1/2}{k}t^k c^{-k}
\)
and summing over \(t=a_i,b_i\), we obtain
\[
G(c)
=c^{1/2}\sum_{k\ge 0}\binom{1/2}{k}c^{-k}\Bigl(\sum_i a_i^k-\sum_i b_i^k\Bigr).
\]
By power sum matching, all \(k\le Q\) terms vanish, while the \(k=q_0\) term coefficient is non-zero, so \(G(c)=\Theta(c^{1/2-q_0})\);
thus for sufficiently large \(c\) we have \(G(c)\ne 0\), i.e., \(\sum_i \sqrt{a_i+c}\ne \sum_i \sqrt{b_i+c}\).
\par
Fix such sufficiently large \(c\), and let the two families of multiplicity multisets be \(A^{(c)}\), \(B^{(c)}\) respectively (further same-scale integer scaling does not change the equal-power-sum relation).
For any sufficiently large \(m\), add the same number of \(1\)s on both sides to fill the total sum to \(2^m\), and we obtain the required integer-valued spectra \(d_m,\widetilde d_m\) (without changing the \(q\le Q\) power sum equality, and without changing the half-order power sum difference).
After normalizing them to \(w_m,\widetilde w_m\), the integer-order collision moments are identical while \(A_{1/2}\) differs.
By the endpoint expansion of Theorem \ref{thm:pom-diagonal-rate-small-distortion-hellinger-endpoint}, the leading-order constant \(A_{1/2}(w)^2-1\) controls the slope term as \(\delta\downarrow 0\), so this endpoint behavior cannot be uniquely determined by integer moments of any fixed \(Q\) as \(m\to\infty\).
\end{proof}

